\chapter{Debate Format \& Moderator Guide} \index{Debate Protocol}


\marginnote{\newthought{Core Rules for Debators}

\begin{itemize}
      \item \textbf{Stay in role.} You have been assigned a position; argue
              it, even if you disagree personally. The purpose is to stress-test
              positions, not to locate the comfortable middle. 
              Thus, you may want to prepare for both sides of the argument in advance.

      \item \textbf{Handshake, then speak.} Every chair entry begins with a
              handshake. The ritual is not optional. It signals that the
              argument is about ideas, not the person.

      \item \textbf{Attack the argument, not the speaker.} Challenge the
              reasoning; do not make it personal.

      \item \textbf{Cite the text or case.} If you reference the assigned
              text or the session stimulus, name it and state what it shows.
              Vague gestures at ``the story'' are not arguments.

      \item \textbf{Concede precisely, then pivot.} If the opponent has
              made a point you cannot rebut, say so clearly and explain why
              your position survives it. A clean concession with a pivot is a
              stronger move than evasion.

      \item \textbf{Recognise quality across the line.} Knock your knuckles
              on the chair when the opposing or your speaker earns it. A debate that
              acknowledges good arguments is a better debate.

      \item \textbf{Use the pass.} Stepping down voluntarily when you are
              out of arguments is better than stalling. It gives your team a
              fresh voice. It is not a defeat.

      \item \textbf{Red-flag sparingly.} A red flag is a public vote of
              no-confidence in your current speaker. Reserve it for when the
              argument is being lost, not when you disagree with the phrasing.
              Remember: a red-flagged speaker may not return until every other
              team member has held the chair.

      \item \textbf{Team: confer quietly.} Members not in the chair may
              confer quietly with each other but may not coach their speaker
              audibly. Take notes; prepare to take the chair.
\end{itemize}}

\noindent Each session debate follows the same three-phase structure. The format is
designed to produce genuine argument, not consensus, but the productive
collision of well-prepared positions, which also do not necessarily need to reflect your own. This guide is for both the moderator
and debators.

\bigskip

\noindent\begin{tabular}{@{}p{3.5cm}p{1.8cm}p{4.5cm}@{}}
    \textbf{Phase} & \textbf{Duration} & \textbf{Activity} \\[4pt]
    1: Primer     & 5 min & Introduction and Opening Question \\
    2: Surrounded & 35 min & Surrounded debate with dynamic chair exchange \\
    3: Wrap-up    & 10 min & Key takeaways and session debrief \\
\end{tabular}

% ---------------------------------------------------------------
\section*{The Surrounded Format}

\noindent The debate is contested by two teams: the \textbf{Proposition team}
(arguing for a motion) and the \textbf{Opposition team} (arguing against it).
Each team has \textbf{3--4 members}. Two chairs are placed facing each other at
the centre of the room. One member from each team (the current speaker)
occupies the centre chair at all times. The rest of their team sits in a
semi-arc \emph{directly behind} their speaker, facing the opposing pair. The
``surrounded'' geometry is thus two opposing benches facing each other across
the two centre chairs.

\medskip
\noindent The debate is conversational, not a sequence of speeches. Speakers
address each other directly. There is no pre-set speaking order after the
opening; the exchange follows the argument.

\section*{Chair Exchange Protocol}

\marginnote{\newthought{The Moderator}

One student serves as moderator per session, rotating each week. The
moderator does not argue a position and sits at the front with a clear
sightline to both teams. Their responsibilities are:

\begin{itemize}
      \item Arrange the room before the session: two chairs facing each
              other at the centre, each team in a row or arc directly behind
              their chair.
      \item Open the session, read the motion, assign teams and sides.
      \item Judge red-flag calls: confirm a simple majority of the bench
              has raised their signal, then call the exchange.
      \item Track the \textbf{rotation state} per team: record who has held
              the chair, so the rotation rule can be enforced accurately.
      \item Logs arguments around \textbf{knocking moments} and \textbf{concessions} as they occur;
              announce each briefly so the room is aware.
      \item Is allowed to initiate concessions and knock, but only in the spirit of maintaining a fair and
              high-quality debate, not to influence the outcome.
      \item Enforce the 60-second minimum seat time and the 60/30-second
              post-knock and post-concession protection windows.
      \item Enforce the handshake ritual at every chair exchange.
      \item Close the debate with a two-sentence neutral summary, noting any
              concessions that shifted the terms of the argument.
      \item Hand over to the instructor for the wrap-up.
\end{itemize}}

\noindent Two mechanisms move speakers in and out of the chairs.

\bigskip

\noindent\textbf{\enspace Voluntary Pass}

\noindent A speaker may tag themselves out at any point by saying \emph{``I
pass.''} No team vote is required. The next available team member takes the
chair. The outgoing and incoming speakers \textbf{shake hands} at the chair.
\emph{No handshake with the opponent;} this is a team substitution, not a
challenge. A voluntary pass does \emph{not} count against the rotation
requirement: the passing speaker may return as soon as a teammate is willing
to swap back.

\bigskip

\noindent\textbf{\enspace Red Flag (forced replacement)}

\noindent A team may replace their current speaker after a \textbf{minimum seat
time of 60 seconds}. To red-flag:

\begin{enumerate}
    \item A simple majority of the team members \emph{not in the chair} raise
          a visible signal simultaneously (a card, a raised hand; agree on
          the signal before the session starts).
    \item The moderator confirms the majority and calls the exchange.
    \item The outgoing and incoming speakers \textbf{shake hands} at the chair.
          The incoming speaker then \textbf{shakes hands with the opponent}
          before speaking.
    \item \textbf{Rotation rule:} a red-flagged speaker may not return to the
          chair until every other team member has held it at least once.
          Full rotation restores eligibility.
\end{enumerate}





% ---------------------------------------------------------------
\section*{Debate Quality Mechanisms}

\noindent Two additional mechanisms are available to all participants. They are
not used to change the chair; they are used to regulate the \emph{quality}
of the exchange, keep it constructive, and make intellectual honesty visible.

\bigskip

\noindent\textbf{Knocking (cross-team acknowledgment)}

\noindent When members of either team believe a speaker (including the
opponent) has made an exceptionally strong argument, they acknowledge it by
\textbf{knocking their knuckles on the armrest or seat of their chair}. The
knock is the academic equivalent of applause: brief, audible, unmistakable.

\begin{enumerate}
    \item Any number of participants may knock spontaneously; no majority
          is required. The knock is individual, not coordinated.
    \item The moderator notes the moment but does not interrupt play.
          Knocking is logged informally as a marker of quality.
    \item Play continues immediately. No chair change occurs. Knocking
          carries no mechanical penalty or reward; it is a public, embodied
          signal of respect for a well-made point.
\end{enumerate}

\noindent\textit{Purpose:} Builds a shared culture of recognising good
argument. The physical gesture is harder to fake than raising a card and
creates an immediate, visceral signal that something important was said.

\bigskip

\noindent\textbf{Point of Concession (speaker-initiated)}

\noindent A speaker in the chair may formally acknowledge that the opponent
has made a point they cannot fully rebut. This is not a sign of weakness;
it is the highest-quality move in academic debate.

\begin{enumerate}
    \item The speaker raises an open hand (or a white card) and says:
          \emph{``I concede: [state the opponent's point clearly and
          precisely].''} Vague concessions are not valid.
    \item The speaker must then pivot: \emph{``\ldots\ my position survives
          because [explanation].''} A concession without a pivot is a
          capitulation and results in a speaker swap; the pivot is mandatory.
    \item The moderator logs the concession and announces it briefly:
          \emph{``Concession noted.''}
    \item The conceding speaker is \textbf{protected from a red flag for
          30 seconds} after the concession. This makes intellectual honesty
          safe to demonstrate.
    \item The opponent \emph{may not repeat or pile on the conceded point}
          for the remainder of that speaking turn. The concession closes
          that line of argument; the debate must move forward.
\end{enumerate}

\noindent\textit{Purpose:} Makes weaknesses transparent rather than papered
over. Rewards calibrated, honest reasoning. Keeps the debate from
cycling on already-settled sub-questions.


% ---------------------------------------------------------------


% ---------------------------------------------------------------
\section*{Moderator Opening Script}

\medskip
\noindent\textit{This is a guide, not a script to read verbatim.
Adapt it to the room.}

\medskip
\noindent \textbf{1.\ Convene:}

\noindent The motion before us today is: [read the motion]. I am the moderator.
I will manage chair exchanges, log concessions, and
keep time. Please take your team seats.

\noindent \textbf{2.\ Assign teams and initial chairs:}

\noindent Team assignments are as follows (if possible done randomly): [read names and roles]. The
Proposition team is [names]; the Opposition team is [names]. The
Proposition side opens with [Speaker A]; the Opposition side opens
with [Speaker B]. Take two minutes now to review your opening
position and agree on your red-flag signal.

\noindent \textbf{3.\ Read the stimulus and open:}

\noindent Before we begin, I will read the session stimulus. [Read the
stimulus from the Session Stimulus section.]

\medskip
\noindent [Speaker A] and [Speaker B]: please take the chairs, face each
other, and shake hands. There are no opening statements; begin
directly. The debate is open.

% ---------------------------------------------------------------
\section*{Moderator Closing Script}

\noindent The exchange is closed. I will now summarise the main fault lines.
[Give a two-sentence neutral summary, noting any concessions that
shifted the ground of the argument.] I hand over to the instructor
for the wrap-up.



